% Clase 22 — solución gráfica del circuito con imán: intersección de la recta
% desmagnetizante B = -mu0 (L_AB/L_CD) H, impuesta por el circuito, con la curva
% de histéresis del imán (cualitativa). Hay dos soluciones simétricas porque no
% se dio la orientación del imán.
\begin{tikzpicture}
  \begin{axis}[
    emfig, width=9.5cm, height=7cm,
    axis lines=middle,
    xlabel style={anchor=north west, font=\small},
    ylabel style={anchor=south east, font=\small},
    xmin=-4.2, xmax=4.2, ymin=-1.65, ymax=1.65,
    xtick=\empty, ytick=\empty, grid=none, clip=false,
    xlabel={$H_m$}, ylabel={$B_m$},
  ]
    \addplot[figblue, domain=-4:4, samples=120] {1.3*tanh((x+1)/0.9)};
    \addplot[figblue, domain=-4:4, samples=120] {1.3*tanh((x-1)/0.9)};
    \addplot[figred, domain=-2.4:2.4, samples=2] {-0.6*x};
    % intersecciones (resueltas por biseccion: H = -0.6989, B = 0.4194)
    \fill[figred] (axis cs:-0.6989,0.4194) circle (2.2pt);
    \fill[figred] (axis cs:0.6989,-0.4194) circle (2.2pt);
    \node[font=\footnotesize, text=figred, anchor=north west] at (axis cs:2.45,-1.05)
      {$B_m = -\mu_0\dfrac{L_{AB}}{L_{CD}}\,H_m$};
  \end{axis}
\end{tikzpicture}
